\section{Controlled and two-qubit primitive gates}

\GateHeading{Controlled NOT}{CNOT}{\ControlledSegment{$\oplus$}}

\[
\ket{c,t}\longmapsto\ket{c,t\oplus c}.
\]
CNOT leaves the control unchanged. When the control is zero, the target is
unchanged; when the control is one, X acts on the target. Because the map is a
permutation of basis states, it is reversible and self-inverse. A
superposition control can entangle the wires, so ``classical XOR'' describes
only its basis-state action.
\begin{lstlisting}
CNOT -I Control -O Target
\end{lstlisting}
\begin{center}
\begin{tabular}{cc|cc}\toprule $c$&$t$&$c'$&$t'$\\\midrule
0&0&0&0\\0&1&0&1\\1&0&1&1\\1&1&1&0\\\bottomrule
\end{tabular}
\end{center}
CNOT is a reversible embedding of XOR because it preserves the control.
In the workbench, the control comes from \textbf{Control A} and the target
from \textbf{Target particle}; \textbf{Control B} is ignored.

\GateHeading{Doubly controlled NOT / Toffoli}{CCNOT}{\DoubleControlledSegment}

\[
\ket{a,b,t}\longmapsto\ket{a,b,t\oplus(a\land b)}.
\]
CCNOT, also called Toffoli, applies X only when both controls are one. It is a
reversible universal building block for classical computation embedded in a
quantum circuit. The target's previous value matters: zero computes AND, while
one computes NAND. The gate is self-inverse.
\begin{lstlisting}
CCNOT -I A B -O Target
\end{lstlisting}
\begin{center}
\begin{tabular}{ccc|ccc}\toprule $a$&$b$&$t$&$a'$&$b'$&$t'$\\\midrule
0&0&0&0&0&0\\0&0&1&0&0&1\\
0&1&0&0&1&0\\0&1&1&0&1&1\\
1&0&0&1&0&0\\1&0&1&1&0&1\\
1&1&0&1&1&1\\1&1&1&1&1&0\\\bottomrule
\end{tabular}
\end{center}
With a target initialized to zero, Toffoli computes classical AND reversibly.

\begin{beginnerbox}[Controls are preserved, the target is the output]
In the table, the $a'$ and $b'$ columns always equal $a$ and $b$: the two
controls are only read. Only $t$ can change, and it flips exactly in the two
rows where $a=b=1$. In the \textbf{Interactive workbench}, choose
\textbf{CCNOT} in \textbf{Gate}, the two inputs in \textbf{Control A} and
\textbf{Control B}, and the output wire in \textbf{Target particle}. The
\emph{CCNOT demo} starter circuit is exactly this gate with both controls
prepared as 1.
\end{beginnerbox}

\GateHeading{Controlled Z}{CZ}{\ControlledSegment{Z}}

\[
\mathrm{CZ}\ket{c,t}=(-1)^{ct}\ket{c,t},\qquad
\mathrm{CZ}=\operatorname{diag}(1,1,1,-1).
\]
CZ adds a minus sign only to the $\ket{11}$ amplitude. No computational-basis
bit changes, but superpositions can acquire observable relative phase. CZ is
symmetric between its two wires, even though QPU syntax writes one as the
input control and one as the output target. It is self-inverse.
\begin{lstlisting}
CZ -I Control -O Target
\end{lstlisting}
\begin{center}
\begin{tabular}{cc|cc}\toprule Input & Measured output & Amplitude multiplier & State\\\midrule
00&00&$+1$&$\ket{00}$\\01&01&$+1$&$\ket{01}$\\
10&10&$+1$&$\ket{10}$\\11&11&$-1$&$-\ket{11}$\\\bottomrule
\end{tabular}
\end{center}
The classical columns are unchanged; the quantum phase is the operation.

\GateHeading{Controlled phase}{CPHASE}{\ControlledSegment{$P(\theta)$}}

\[
\ket{11}\longmapsto e^{i\theta}\ket{11},
\]
while $\ket{00}$, $\ket{01}$, and $\ket{10}$ are unchanged. CZ is the special
case $\theta=\pi$. CPHASE unlocks compact Quantum Fourier Transform and
phase-estimation demos.
\begin{lstlisting}
CPHASE=pi/2 -I Control -O Target
CPHASE=90d -I Control -O Target
\end{lstlisting}

\GateHeading{Controlled Y}{CY}{\ControlledSegment{Y}}

\[
\ket{0,t}\mapsto\ket{0,t},\quad
\ket{1,0}\mapsto i\ket{1,1},\quad
\ket{1,1}\mapsto-i\ket{1,0}.
\]
CY applies the Pauli Y rotation to the target only when the control is one. Its
basis-bit pattern resembles CNOT, but the $\pm i$ factors distinguish their
behavior on superpositions. Since $Y^2=I$, CY is self-inverse.
\begin{lstlisting}
CY -I Control -O Target
\end{lstlisting}
\begin{center}
\begin{tabular}{cc|cc|c}\toprule $c$&$t$&$c'$&$t'$&Quantum phase\\\midrule
0&0&0&0&$+1$\\0&1&0&1&$+1$\\
1&0&1&1&$i$\\1&1&1&0&$-i$\\\bottomrule
\end{tabular}
\end{center}

\GateHeading{Swap}{SWAP}{\SwapSegment}

\[
\mathrm{SWAP}\ket{a,b}=\ket{b,a}.
\]
SWAP exchanges the complete quantum states of two wires, including amplitudes
and any correlations with the rest of the register. It does not copy either
state. Three CNOTs can decompose SWAP, but the QPU system provides it as one
primitive instruction. Applying SWAP twice restores the original ordering.
\begin{lstlisting}
SWAP -I A B -O A B
\end{lstlisting}
The two output tokens are mandatory and must resolve to the same two wires, in
the same order, as the inputs.
\begin{center}
\begin{tabular}{cc|cc}\toprule $a$&$b$&$a'$&$b'$\\\midrule
0&0&0&0\\0&1&1&0\\1&0&0&1\\1&1&1&1\\\bottomrule
\end{tabular}
\end{center}

